\textbf{结论名称：} 二面角的求法三：射影面积法

\textbf{适用条件：} 已知平面多边形（或一般平面图形）在两相交平面之一上的正投影面积，用于求两平面所成二面角。当前图形无折叠重叠，投影为垂直投影。

\textbf{变量说明：} $\theta$ 为两平面所成二面角的平面角（$0<\theta<\pi$）；$S_{\text{原}}$ 为原平面图形的面积；$S_{\text{投影}}$ 为该图形在另一平面上的正投影面积。

\textbf{核心公式：} 
\[
\boxed{|\cos\theta| = \frac{S_{\text{投影}}}{S_{\text{原}}}}
\]
当可以确定二面角为锐角时，$\cos\theta = \dfrac{S_{\text{投影}}}{S_{\text{原}}}$；若为钝角则取负号。

\textbf{使用场景：} 不易直接作出二面角的平面角时，可转化为面积比计算。常见于棱柱的截面与底面、棱锥的侧面与底面所成二面角等问题。

\textbf{简单例子：} 矩形 $ABCD$ 面积为 $4$，它在平面 $\beta$ 上的正投影面积为 $2\sqrt{3}$，且矩形所在平面与 $\beta$ 所成二面角为锐角 $\theta$。则
\[
\cos\theta = \frac{2\sqrt{3}}{4} = \frac{\sqrt{3}}{2},\quad \text{故 } \theta = 30^\circ.
\]